% Compile with: xelatex uct_math_en.tex
\documentclass[conference]{IEEEtran}
\usepackage{fontspec}
\usepackage[english]{babel}
\usepackage{amsmath,amssymb,amsthm,mathtools}
\usepackage{hyperref,algorithm,algpseudocode,enumitem}
\setmainfont{Liberation Serif}
\hypersetup{colorlinks=true,linkcolor=blue,citecolor=blue,urlcolor=blue}
\newtheorem{definition}{Definition}[section]
\newtheorem{proposition}[definition]{Proposition}
\newtheorem{theorem}[definition]{Theorem}
\newtheorem{remark}[definition]{Remark}
\newcommand{\E}{\mathbb{E}}
\newcommand{\Prob}{\mathbb{P}}
\title{Mathematical Foundations and Analysis of UCT\\
\large An Application to Monte Carlo Tree Search for Tic-Tac-Toe/Gomoku}
\author{\IEEEauthorblockN{Student Author}
\IEEEauthorblockA{Department of Artificial Intelligence\\
University or Institution}}
\begin{document}
\maketitle
\begin{abstract}
Upper Confidence Bound applied to Trees (UCT) is a node-selection rule for Monte
Carlo Tree Search (MCTS), derived from the multi-armed bandit problem. This paper
formulates finite adversarial games, derives UCB1 from a concentration bound,
explains the exploration constant, and analyzes Selection, Expansion, Simulation,
and Backpropagation. It maps the mathematics to the existing five-in-a-row code
and gives a bonus formulation for 3x3 Tic-Tac-Toe. Theoretical guarantees are
separated from empirical claims; no unmeasured numbers are reported.
\end{abstract}
\begin{IEEEkeywords}
MCTS, UCT, UCB1, multi-armed bandits, adversarial games, Tic-Tac-Toe, Gomoku
\end{IEEEkeywords}

\section{Introduction}
Monte Carlo Tree Search incrementally constructs a search tree using simulations
instead of exhaustively expanding the game tree. UCT is a standard Selection rule:
it prefers actions with high empirical reward while reserving computation for
under-sampled actions \cite{kocsis2006}.

\paragraph{Note.} ``MCTS combined with UCT'' is most precisely understood as MCTS
using UCT during Selection. The rollout policy, terminal test, and reward convention
are separate design choices.

\section{Mathematical game model}
\begin{definition}[Finite adversarial game]
Represent a game by
\[
\mathcal G=(\mathcal S,\mathcal A,T,s_0,\mathcal Z,r),
\]
where $\mathcal S$ is the state space, $\mathcal A(s)$ is the legal-action set,
$T(s,a)$ is the successor state, $s_0$ is the initial state, $\mathcal Z$ is the
terminal set, and $r\in[0,1]$ is a reward from one player's perspective.
\end{definition}

For Tic-Tac-Toe, a state contains a 3x3 board and the player to move. For Gomoku,
the board may be $n\times n$ and a player wins after forming at least five
consecutive marks horizontally, vertically, or diagonally. Legal actions are empty
cells; terminal states are wins or full boards.

\paragraph{Note.} The transition function must enforce legality. If a rollout can
play an occupied cell, its samples no longer represent the intended game.

\section{From multi-armed bandits to UCB1}
\subsection{The bandit problem}
Consider $K$ arms. Pulling arm $i$ gives a random reward $X_i\in[0,1]$ with mean
$\mu_i$. After $t$ rounds, define
\[
N_i(t)=\sum_{\tau=1}^{t}\mathbf 1\{I_\tau=i\},\qquad
\widehat\mu_i(t)=\frac{1}{N_i(t)}\sum_{\tau:I_\tau=i}X_\tau.
\]
Let $\mu^*=\max_i\mu_i$ and $\Delta_i=\mu^*-\mu_i$. Expected regret is
\[
R_T=\E\!\left[\sum_{t=1}^{T}(\mu^*-\mu_{I_t})\right]
=\sum_{i=1}^{K}\Delta_i\E[N_i(T)].
\]

\subsection{Confidence bounds and UCB1}
For independent variables in $[0,1]$, Hoeffding's inequality gives
\[
\Prob(|\widehat\mu_i-\mu_i|\geq\varepsilon)
\leq 2\exp(-2N_i\varepsilon^2).
\]
Using $b_i(t)=\sqrt{2\ln(t)/N_i(t)}$ gives the UCB1 index
\[
U_i(t)=\widehat\mu_i(t)+\sqrt{\frac{2\ln t}{N_i(t)}}.
\]
The algorithm selects an arm maximizing this index. The first term exploits; the
second explores. UCB1 has logarithmic regret bounds in $T$ under standard bandit
assumptions \cite{auer2002}.

\begin{proposition}[Interpretation of UCB]
When $N_i(t)$ is small, the confidence radius is large and arm $i$ is likely to be
sampled again. As $N_i(t)$ grows, the radius decreases as $N_i(t)^{-1/2}$ and
the decision increasingly reflects the empirical mean.
\end{proposition}

\paragraph{Note.} UCT commonly uses
\[
\widehat\mu_j+c\sqrt{\frac{\ln N(s)}{N(s,j)}},
\]
with $c>0$. $c=\sqrt2$ is a common UCB1 normalization; $c=1.4$ is only a
practical starting point, not a universal optimum.

\section{UCT on a game tree}
\begin{definition}[Search node]
A node $s$ stores its visit count $N(s)$, accumulated reward $W(s)$, parent,
expanded children, and untried actions. Its empirical value is
\[
Q(s)=W(s)/N(s)\quad\text{when }N(s)>0.
\]
\end{definition}

At state $s$, for a child corresponding to action $a$, UCT selects
\[
a^*\in\arg\max_a\left[Q(s,a)+c\sqrt{\frac{\ln N(s)}{N(s,a)}}\right].
\tag{1}
\]
If $N(s,a)=0$, the index is set to $+\infty$, so every new child is sampled.

\subsection{The four MCTS phases}
\begin{enumerate}[label=\arabic*.]
\item \textbf{Selection:} descend from the root using (1) to a terminal node or
a node with an untried action.
\item \textbf{Expansion:} choose an untried action, create a child, and apply it.
\item \textbf{Simulation:} roll out to termination using a random or heuristic policy.
\item \textbf{Backpropagation:} propagate the outcome back along the path and update
$N$ and $W$.
\end{enumerate}

\begin{algorithm}[h]
\caption{MCTS--UCT for a two-player game}
\begin{algorithmic}[1]
\Require root state $s_0$, budget $M$, exploration constant $c$
\State create root node $v\gets s_0$
\For{$k=1$ to $M$}
  \State $u\gets v$
  \While{$u$ is fully expanded and nonterminal}
    \State $u\gets\arg\max_x\left(Q(x)+c\sqrt{\ln N(u)/N(x)}\right)$
  \EndWhile
  \If{$u$ is nonterminal and has an untried action}
    \State create a child from an untried action; set $u$ to the new node
  \EndIf
  \State $z\gets\mathrm{Rollout}(u)$; backpropagate $z$ from $u$ to $v$
\EndFor
\State \Return the root child with the largest visit count
\end{algorithmic}
\end{algorithm}

\paragraph{Note.} Choosing the most visited child is not the same as choosing the
largest $Q$. Visit count is often more stable when a high mean is based on very few
samples; it is also the rule used by the current code.

\section{Probability and asymptotic analysis}
At a fixed state, if each action is treated as an arm and rollout rewards are
independent and bounded in $[0,1]$, UCB1 intuition applies locally. In a real tree,
allocation is adaptive and policies below nodes change, so samples are not fully
independent.

\begin{theorem}[Asymptotic consistency of UCT]
Under the standard assumptions for UCT on a finite tree, bounded rewards, and valid
rollouts, as the number of simulations tends to infinity, the probability that UCT
selects a suboptimal root action tends to zero \cite{kocsis2006}.
\end{theorem}

This is not a guarantee of a win with a small finite budget. Deep trees, rare
rewards, trap actions, biased rollouts, and time limits may make finite-budget
performance far from the limit.

For $U=Q+c\sqrt{\ln N(s)/N(s,a)}$, $c=0$ favors exploitation, large $c$ favors
exploration, and intermediate values require empirical tuning.

\paragraph{Note.} Do not claim that one $c$ is best from one game. Use multiple
seeds, swap the first player, and report wins, draws, losses, time, nodes, and the
opponent definition.

\section{Mapping to the existing Gomoku implementation}
In \texttt{tic\_tac\_toe\_ai.py}, \texttt{MCTSNode} stores \texttt{visits} and
\texttt{wins}. \texttt{uct(exploration)} computes
\[
W/N+c\sqrt{\ln N_{\mathrm{parent}}/N},
\]
returning infinity for an unvisited node. \texttt{best\_mcts\_move} performs
Selection, Expansion, calls \texttt{rollout}, and performs Backpropagation for
the number of iterations given by \texttt{--iterations}.

The reward is measured from the player who just moved:
\[
z=\begin{cases}1,&\text{win},\\0.5,&\text{draw},\\0,&\text{loss}.
\end{cases}
\]
This convention must be consistent in both UCT and Backpropagation. A heuristic
rollout that prefers immediate wins or blocking moves changes the sampling
distribution; it should not be described as unbiased random play.

\section{Bonus: 3x3 Tic-Tac-Toe}
For Tic-Tac-Toe, $\mathcal A(s)$ is the set of empty cells, $T$ places the mark of
the player to move, and terminal rewards are $1,0,0.5$. Each iteration expands an
empty cell, rolls out until three-in-a-row or a full board, updates the path, and
selects the root action with the largest visit count. Tic-Tac-Toe is useful because
minimax can solve it exactly, allowing UCT to be evaluated against an optimal player
rather than by counting one observed win.

\section{Suggested experimental protocol}
Do not insert invented measurements. Compare $c\in\{0.5,1,\sqrt2,2\}$ and
$M\in\{100,500,1000,5000\}$, use at least 30 seeds per configuration, compare
random and heuristic rollouts, and report win/draw/loss rates, move time, node
count, $N(s_0)$, board size, first player, win rule, and opponent.

\section{Conclusion}
UCT brings confidence bounds to a game tree: $Q$ exploits and the square-root term
explores. ``Converges with infinitely many simulations'' does not mean ``always
wins with a finite budget.'' Performance in Gomoku and Tic-Tac-Toe depends on
legal move generation, rollout policy, rewards, $c$, and the simulation budget.

\begin{thebibliography}{9}
\bibitem{auer2002} P. Auer, N. Cesa-Bianchi, and P. Fischer, ``Finite-time
Analysis of the Multiarmed Bandit Problem,'' \emph{Machine Learning}, 47,
235--256, 2002. \href{https://doi.org/10.1023/A:1013689704352}{doi}.
\bibitem{kocsis2006} L. Kocsis and C. Szepesvári, ``Bandit Based Monte-Carlo
Planning,'' ECML, LNCS 4212, 282--293, 2006.
\href{https://doi.org/10.1007/11871842_29}{doi}.
\bibitem{browne2012} C. B. Browne et al., ``A Survey of Monte Carlo Tree Search
Methods,'' \emph{IEEE TCIAIG}, 4(1), 1--43, 2012.
\href{https://doi.org/10.1109/TCIAIG.2012.2186810}{doi}.
\bibitem{sutton2018} R. S. Sutton and A. G. Barto, \emph{Reinforcement Learning:
An Introduction}, 2nd ed., MIT Press, 2018.
\href{http://incompleteideas.net/book/the-book-2nd.html}{online edition}.
\end{thebibliography}
\end{document}
